% pipe_opt2 (2026-11-02): OPTION 2 "accent bars". Layout = mock_A_full_v2_8pt1col. No large colored fills: clause slabs are white
% with a thick clause-colored left bar; stage titles are dark tabs; sources of truth are green icons on a light-green rail.
% SINGLE-COLUMN 8pt-minimum version of mock_A_full_v2 (2026-11-02): every \scriptsize -> \footnotesize, macros from mock_macros_8pt.tex, spacing squeezed.
% Direction A, full text, v2 (2026-11-02): + 3-pronged DSL under Stage-1 clause 3; Stage-2 passes as a
% clear vertical chain of numbered chevron badges (v1 stripes were too subtle). v1 = mock_A_full.tex.
% Direction A, full text ("A_full"): A's visual language carrying ALL the wording of the main figure
% (paper_figures/drafts/pipeline_vertical.tex as of 2026-11-02 01:49).
% Changes from mock_A to make room for the text: every stage flows top-down (clauses stacked
% full-width, the solver BELOW them), so the data spine runs straight down the left with no
% return elbows; the Stage-2 passes are full-width downward chevrons ("stripes"), one per line.
\documentclass[tikz,border=0pt]{standalone}
\usepackage{newtxtext,newtxmath}
\usepackage{graphicx}
% pipe_opt_macros.tex (2026-11-02): copy of mock_macros_8pt.tex for pipe_opt1/2/3; \sot takes its fill/draw from \sotfill/\sotdraw, + \chip.
% 8pt-minimum variant of mock_macros.tex (2026-11-02): every text style is \footnotesize (8pt);
% icon glyphs that were \tiny text ({ } and </>) are drawn as paths instead.
% Shared shape vocabulary for the big-figure design mocks (BIGFIG_SHAPE_DESIGN.md).
% Every tikzpicture that uses these must set [x=1pt,y=1pt]: all numbers below are pt.
\usetikzlibrary{calc,fit,shapes.symbols,shapes.misc,shapes.geometric,decorations.pathmorphing,arrows.meta}
\def\defshade{black!12}   % the one gray: defense-specific
\def\sotfill{\defshade}\def\sotdraw{black}   % source-of-truth document icon colors (redefine per option)
\tikzset{
  >=latex,
  every node/.style={font=\footnotesize, inner sep=0pt,
     execute at begin node={\hyphenpenalty=10000\exhyphenpenalty=10000\relax}},
  q/.style={font=\footnotesize\itshape},             % qualifier: multiplicity, mode, arrow label
  sm/.style={font=\footnotesize},                    % small roman label
  blob/.style={draw, rounded corners=7pt, line width=0.6pt},          % stage container
  btitle/.style={fill=white, inner xsep=2pt, inner ysep=1pt, anchor=west}, % title sitting on the blob border
  slab/.style={draw, rounded corners=2.5pt, minimum height=10.5pt, inner xsep=3pt, align=left, anchor=north west},
  ghost/.style={slab, dash pattern=on 1.4pt off 1.1pt},               % "same as Stage 1"
  pill/.style={draw, rounded corners=4.5pt, minimum height=10.5pt, inner xsep=3.5pt, align=center},  % formula pill
  tool/.style={draw, rounded rectangle, minimum height=17pt, inner xsep=5pt, align=center, line width=0.6pt},  % solver / engine
  chev/.style={draw, signal, signal to=east, signal from=west, signal pointer angle=125,
               minimum height=19pt, align=center, inner xsep=1pt, font=\footnotesize},
  chevL/.style={draw, signal, signal to=west, signal from=east, signal pointer angle=125,
               minimum height=19pt, align=center, inner xsep=1pt, font=\footnotesize},
  chevD/.style={draw, signal, signal to=south, signal from=north, signal pointer angle=140,
               minimum width=15pt, minimum height=10pt, align=center, inner sep=0pt, font=\footnotesize\bfseries},
  dspec/.style={fill=\defshade},
  dot/.style={circle, fill, inner sep=0.75pt},
  station/.style={circle, draw, fill=white, minimum size=8.5pt, inner sep=0pt, font=\footnotesize\bfseries},
  box/.style={draw, rounded corners=3pt, align=center, inner xsep=3pt, inner ysep=2.2pt},
}

% ---------- glyphs drawn inside a card centered at the scope origin (card 11 x 14) ----------
\newcommand{\gGraph}{% execution graph: 3 instruction nodes in spo + one rf arc (matches fig a)
  \draw[line width=0.3pt] (-1.8,3.6) -- (-1.8,-3.6);
  \fill (-1.8,3.6) circle (0.95) (-1.8,0) circle (0.95) (-1.8,-3.6) circle (0.95);
  \draw[line width=0.35pt, -{Latex[length=1.5pt,width=1.5pt]}] (-1.0,3.8) .. controls (3.6,3.2) and (3.6,-3.2) .. (-0.9,-3.5);}
\newcommand{\gAsm}{% program text
  \foreach \y/\l in {4/5.2,2/3.8,0/5.2,-2/2.8,-4/4.4} \draw[line width=0.35pt] (-3,\y) -- ++(\l,0);}
\newcommand{\gAnn}{% annotation: a pair of curly braces, drawn (no text)
  \draw[line width=0.4pt, line join=round] (-1.2,3.8) .. controls (-2.4,3.8) and (-2.0,1.0) .. (-2.4,0.3) -- (-3.1,0) -- (-2.4,-0.3) .. controls (-2.0,-1.0) and (-2.4,-3.8) .. (-1.2,-3.8);
  \draw[line width=0.4pt, line join=round] (1.2,3.8) .. controls (2.4,3.8) and (2.0,1.0) .. (2.4,0.3) -- (3.1,0) -- (2.4,-0.3) .. controls (2.0,-1.0) and (2.4,-3.8) .. (1.2,-3.8);}
\newcommand{\gCode}{% source code: </> drawn (no text)
  \draw[line width=0.45pt, line join=miter] (-1.9,0.6) -- (-4.0,-1.2) -- (-1.9,-3.0) (1.9,0.6) -- (4.0,-1.2) -- (1.9,-3.0) (0.9,1.0) -- (-0.9,-3.4);}
\newcommand{\gTrace}{% pipeline trace: staggered stage bars (O3PipeView-like)
  \foreach \y/\a/\b in {4/-3.6/0.4,1.6/-2.4/1.6,-0.8/-1.2/2.8,-3.2/0/3.6} \draw[line width=0.9pt] (\a,\y) -- (\b,\y);}

% ---------- card stack: #1 name, #2 (x,y) of the front card's center, #3 number of cards, #4 glyph macro ----------
% back cards sit up-right of the front card; anchors #1-n (top of the whole stack), #1-s, #1-e, #1-w
\def\cw{11}\def\ch{14}\def\co{1.8}
\newcommand{\cardstack}[4]{%
  \begin{scope}[shift={#2}]
    \foreach \k in {#3,...,1} {\pgfmathsetmacro\o{\co*(\k-1)}
      \draw[fill=white, rounded corners=0.9pt] (-\cw/2+\o,-\ch/2+\o) rectangle ++(\cw,\ch);}
    #4
    \pgfmathsetmacro\top{\ch/2+\co*(#3-1)}\pgfmathsetmacro\rgt{\cw/2+\co*(#3-1)}
    \path (\co*#3/2-\co/2,\top) coordinate (#1-n) (0,-\ch/2) coordinate (#1-s)
          (\rgt,\co*#3/2-\co/2) coordinate (#1-e) (-\cw/2,0) coordinate (#1-w);
  \end{scope}}
% stack of program + annotation PAIRS: each card is two half-cards clipped together
\newcommand{\pairstack}[3]{% #1 name, #2 (x,y), #3 n
  \begin{scope}[shift={#2}]
    \foreach \k in {#3,...,1} {\pgfmathsetmacro\o{\co*(\k-1)}
      \draw[fill=white, rounded corners=0.9pt] (-11+\o,-7+\o) rectangle ++(10.4,14);
      \draw[fill=white, rounded corners=0.9pt] (0.6+\o,-7+\o) rectangle ++(10.4,14);
      \draw[line width=0.5pt] (-1.4+\o,4.5+\o) -- (1.4+\o,4.5+\o);}  % clip joining the two halves
    \begin{scope}[shift={(-5.8,0)}]\gAsm\end{scope}
    \begin{scope}[shift={(5.8,0)}]\gAnn\end{scope}
    \pgfmathsetmacro\top{7+\co*(#3-1)}
    \path (-5.8,7) coordinate (#1-asm-n) (5.8,7) coordinate (#1-ann-n) (-5.8,-7) coordinate (#1-asm-s) (5.8,-7) coordinate (#1-ann-s)
          (11,0) coordinate (#1-e) (-11,0) coordinate (#1-w) (0,\top) coordinate (#1-n) (0,-7) coordinate (#1-s);
  \end{scope}}

% ---------- defense-specific source of truth: gray folded document; #1 name, #2 (x,y), #3 glyph or empty ----------
\newcommand{\sot}[3]{%
  \begin{scope}[shift={#2}]
    \draw[fill=\sotfill, draw=\sotdraw] (-5.5,-7) -- (5.5,-7) -- (5.5,4) -- (2.5,7) -- (-5.5,7) -- cycle;
    \draw[draw=\sotdraw] (2.5,7) -- (2.5,4) -- (5.5,4);
    \if\relax\detokenize{#3}\relax
      \foreach \y in {1.8,-0.8,-3.4} \draw[line width=0.3pt, draw=\sotdraw] (-3.2,\y) -- (3.2,\y);
    \else #3\fi
    \path (0,7) coordinate (#1-n) (0,-7) coordinate (#1-s) (5.5,0) coordinate (#1-e) (-5.5,0) coordinate (#1-w);
  \end{scope}}

% ---------- bug glyph (#1 position, #2 scale, #3 color) ----------
\newcommand{\buggl}[3]{%
  \begin{scope}[shift={#1}, scale=#2, line width=0.45pt, draw=#3, fill=#3]
    \foreach \y in {1.0,0,-1.0} {\draw (-2.6,\y-0.45) -- (2.6,\y+0.45); \draw (-2.6,\y+0.45) -- (2.6,\y-0.45);}
    \fill (0,-0.2) ellipse (1.35 and 1.9);
    \fill (0,2.15) circle (0.8);
    \draw (-0.35,2.7) -- (-1.2,3.7) (0.35,2.7) -- (1.2,3.7);
  \end{scope}}
% magnifier (#1 center of lens, #2 radius)
\newcommand{\magn}[2]{\draw[line width=0.6pt] #1 circle (#2); \draw[line width=1.1pt] ($#1+(-45:#2)$) -- ++(-45:#2*0.9);}
% legend swatch
\newcommand{\swatch}[1]{\draw[fill=\defshade] #1 rectangle ++(6,6);}

% clause chip: #1 clause key (spec/prot/xmit), #2 text  (use inside a node or as node text)
\newcommand{\chip}[2]{\tikz[baseline=(c.base)]\node[solid, draw=ss#1ink, fill=ss#1fill, rounded corners=1.6pt, inner xsep=1.4pt, inner ysep=0.9pt, line width=0.45pt, font=\footnotesize\bfseries] (c) {\strut#2};}

% SimSpect figure palette (2026-11-03, Anna). Light fills only; all text stays black except red transmitters/bugs.
% Blue = defense-specific / defense sources of truth. Tan = speculation (speculative rows, unresolved window, 3a slab).
% Lavender = protection set (protected state, 3b slab). Light red = leakage function (3c slab, transmitter cells);
% red ink = transmitters, transmissions, real bugs. Load once (paper preamble or a standalone preview) before any body.
% (The earlier stronger per-clause palette is parked in sscolors_clauses.tex.)
\definecolor{sssrcfill}{HTML}{DCE8F4}\definecolor{sssrcbg}{HTML}{DCE8F4}\definecolor{sssrcink}{HTML}{2F5F8F}% defense-specific / sources of truth
\definecolor{ssspecfill}{HTML}{FCEFD6}\definecolor{ssspecink}{gray}{0}% speculation: light tan background
\definecolor{ssprotfill}{HTML}{E4DAF1}\definecolor{ssprotink}{gray}{0}% protection set: light lavender
\definecolor{ssxmitfill}{HTML}{F7E0DD}\definecolor{ssxmitink}{HTML}{B3322C}% leakage: light red fill; red ink for transmitters / bugs

\pgfdeclarelayer{bg}\pgfsetlayers{bg,main}
\def\sotfill{sssrcfill}\def\sotdraw{sssrcink}
\begin{document}
\begin{tikzpicture}[x=1pt,y=1pt]
% ---------------- knobs ----------------
\def\LX{24}      % centre of the sources-of-truth lane
\def\PL{56.3}    % blob left edge
\def\PR{251.7}   % blob right edge
\def\IL{62}      % inner left
\def\SX{80}      % data spine x
\def\SLW{171}    % slab text width (slabs span IL .. IL+177)
\def\CB{243}     % x of the collector bar
\def\BG{32}      % vertical gap between blobs (holds the output strip)
\def\BP{2.5}       % padding below the last element of a blob
\tikzset{
  fslab/.style={slab, text width=\SLW pt},
  fghost/.style={ghost, text width=\SLW pt},
  ftool/.style={draw, rounded rectangle, minimum height=19pt, inner xsep=4pt, align=center, line width=0.6pt, anchor=west},
  passchev/.style={draw, fill=white, signal, signal to=south, signal from=north, signal pointer angle=100,
                 minimum width=14pt, minimum height=13pt, inner sep=0pt, font=\footnotesize\bfseries, line width=0.6pt},
  passtxt/.style={anchor=west, align=left},
  hdr/.style={font=\footnotesize\itshape, anchor=north west},
  lang/.style={font=\footnotesize\itshape, anchor=north west},
  sotl/.style={sm, anchor=north, align=center, text width=47pt, fill=white, inner ysep=1pt},
  capt/.style={anchor=north west, align=left, text width=148pt},
  accent/.style={rounded corners=1.2pt, path picture={\fill[#1] (path picture bounding box.south west) rectangle ($(path picture bounding box.north west)+(2.2pt,0)$);}},
  accent3/.style={rounded corners=1.2pt, path picture={\path (path picture bounding box.south west) coordinate (pw0) (path picture bounding box.north west) coordinate (pw3);
      \fill[ssxmitink] (pw0) rectangle ($(pw0)!1/3!(pw3)+(2.2pt,0)$); \fill[ssprotink] ($(pw0)!1/3!(pw3)$) rectangle ($(pw0)!2/3!(pw3)+(2.2pt,0)$); \fill[ssspecink] ($(pw0)!2/3!(pw3)$) rectangle ($(pw3)+(2.2pt,0)$);}},
  btitle/.style={fill=black!82, text=white, rounded corners=2pt, inner xsep=3pt, inner ysep=1.1pt, anchor=west},
}
\newcommand{\gety}[2]{\path #1; \pgfgetlastxy{\tmpx}{\tmpy}\pgfmathsetmacro#2{\tmpy}}
\coordinate (lxc) at (\LX,0); \coordinate (spx) at (\SX,0); \coordinate (prx) at (\PR,0);
\path (0,-22) -- (252,-22); % fixes the width at exactly 252pt; height follows the content

% ---------------- lane header ----------------
\node[q, anchor=north, align=center, text width=47pt] (lh) at (\LX,-24) {\textbf{Defense\\sources of truth}};

% ================= Stage 1 =================
\def\T{-22}
\node[lang, align=left] at (\IL+4,\T-6.3) {(Alloy relational modeling language):\\[-1.5pt]graph constraints};

\node[fslab] (a1) at (\IL,\T-23.3) {\textbf{1.}~Defense-agnostic execution well-formedness};
\node[fslab] (a2) at ($(a1.south west)+(0,-1.2)$) {\textbf{2.}~Defense-agnostic minimality requirements (only load-bearing components specified)};
\node[fslab, accent=black!50] (a3) at ($(a2.south west)+(0,-1.2)$) {\hspace*{0.6pt}\textbf{3.}~Defense-specific leakage predicate:};
\pgfmathsetmacro\SLWs{\SLW-10}
\node[slab, accent=ssspecink, text width=\SLWs pt] (d1) at ($(a3.south west)+(10,-1.2)$) {\hspace*{0.6pt}\textbf{3a.}~Speculation contract: what is speculative};
\node[slab, accent=ssprotink, text width=\SLWs pt] (d2) at ($(d1.south west)+(0,-1.2)$) {\hspace*{0.6pt}\textbf{3b.}~Protection set: protected state + propagation};
\node[slab, accent=ssxmitink, text width=\SLWs pt] (d3) at ($(d2.south west)+(0,-1.2)$) {\hspace*{0.6pt}\textbf{3c.}~Leakage function: which operands transmit};
\coordinate (tx) at ($(a3.south west)+(6,0)$); \draw (tx) |- (d3.west); \draw (tx |- d1.west) -- (d1.west); \draw (tx |- d2.west) -- (d2.west);
\node[ftool] (A1) at ($(d3.south west)+(-10,-15)$) {Alloy Analyzer\\[-1.5pt]{\footnotesize (Kodkod model finder + SAT backend) + bound}};
\draw[line width=0.6pt] (a1.north -| \CB,0) -- (d3.south -| \CB,0);
\draw[->] (d3.south -| \CB,0) |- (A1.east);

\coordinate (b1) at ($(A1.south)+(0,-\BP)$);
\draw[blob] (\PL,\T) rectangle (b1 -| prx);
\node[btitle] at (\IL+2,\T) {\textbf{Stage 1}\enspace Abstract test generation};
\sot{ds}{(a3.west -| lxc)}{}
\node[sotl] (dsl) at ($(ds-s)+(0,-1.5)$) {Defense specification\\\textit{manual translation}};
\draw[->, sssrcink, line width=0.6pt] (ds-e) -- (a3.west);

% ---- output 1 ----
\gety{(b1)}{\G}
\draw[->] (A1.south -| spx) -- (\SX,\G-5.6);
\cardstack{o1}{(\SX-1.8,\G-15.2)}{3}{\gGraph}
\node[capt] (c1) at (\SX+14,\G-5.5) {Bounded-exhaustive set of defense-specific leaking execution graphs};

% ================= Stage 1.5 =================
\gety{(c1.south)}{\T}\pgfmathsetmacro\T{\T-7.5}
\node[lang, align=left] at (\SX+10,\T-6.3) {(Alloy relational modeling language):\\[-1.5pt]constraints};

\node[fslab] (b4) at (\IL,\T-23.3) {\textbf{+\,4.}~Pin: graph contains target Stage-1 graph};
\draw[->] (o1-s) -- (o1-s |- b4.north);
\node[fghost] (bb1) at ($(b4.south west)+(0,-1.2)$) {\textbf{1.}};
\node[q, anchor=east] at ($(bb1.east)+(-3,0)$) {Same as Stage 1};
\node[fslab] (b2) at ($(bb1.south west)+(0,-1.2)$) {\textbf{2'.}~Loosened minimality requirements $\Rightarrow$ more flexible execution graphs};
\node[fghost, accent3] (bb3) at ($(b2.south west)+(0,-1.2)$) {\hspace*{0.6pt}\textbf{3.}};
\node[q, anchor=east] at ($(bb3.east)+(-3,0)$) {Same as Stage 1};
\node[ftool] (A2) at ($(bb3.south west)+(0,-15)$) {Alloy Analyzer\\[-1.5pt]{\footnotesize (Kodkod model finder + SAT backend) + bound}};
\draw[line width=0.6pt] (b4.north -| \CB,0) -- (bb3.south -| \CB,0);
\draw[->] (bb3.south -| \CB,0) |- (A2.east);

\coordinate (b2b) at ($(A2.south)+(0,-\BP)$);
\draw[blob] (\PL,\T) rectangle (b2b -| prx);
\node[btitle] at (\SX+8,\T) {\textbf{Stage 1.5}\enspace Abstract test enumeration};

% ---- output 2 ----
\gety{(b2b)}{\G}
\draw[->] (A2.south -| spx) -- (\SX,\G-5.6);
\cardstack{o2}{(\SX-2.7,\G-16.4)}{4}{\gGraph}
\node[capt] (c2) at (\SX+14,\G-5.5) {More flexible defense-specific leaking execution graphs\\[-1pt]{\footnotesize\itshape Non-exhaustive; multiple per Stage-1 graph; parallelizable}};

% ================= Stage 2 =================
\gety{(c2.south)}{\T}\pgfmathsetmacro\T{\T-7.5}
\node[lang] at (\SX+10,\T-6.3) {(Python, LLVM): passes, run per graph};

\coordinate (p0) at (\IL+9,\T-16.8);
\foreach \i/\t [count=\k from 0] in {
   1/{Order instructions ($\rightarrow$ PCs); collect state},
   2/{Concretize instructions},
   3/{Concretize speculation primitives (branches, STL)},
   4/{Interleave instruction chains},
   5/{Assign registers and memory; reserve registers},
   6/{SSA + LLVM emission with PC labels},
   7/{Annotations; LLVM $\rightarrow$ asm; PCs $\rightarrow$ addresses}} {
  \path ($(p0)+(0,{-(\i-1)*11.9})$) coordinate (n\i-top);
  \draw[fill=white, line width=0.6pt, line join=miter] ($(n\i-top)+(-7,0)$) -- ($(n\i-top)+(0,-3.6)$) -- ($(n\i-top)+(7,0)$) -- ($(n\i-top)+(7,-10.6)$) -- ($(n\i-top)+(0,-14.2)$) -- ($(n\i-top)+(-7,-10.6)$) -- cycle;
  \coordinate (n\i-c) at ($(n\i-top)+(0,-7.6)$);
  \coordinate (n\i-bot) at ($(n\i-top)+(0,-14.2)$);
  \node[font=\footnotesize\bfseries] at (n\i-c) {\i};
  \node[passtxt] (p\i) at ($(n\i-c)+(10.5,0)$) {\t};
}
\draw[->] (o2-s) |- ($(n1-top)+(0,3)$) -- ($(n1-top)+(0,-3.6)$);
\coordinate (b3) at ($(n7-bot)+(0,-\BP)$);
\draw[blob] (\PL,\T) rectangle (b3 -| prx);
\node[btitle] at (\SX+8,\T) {\textbf{Stage 2}\enspace Concretization};

% ---- output 3 ----
\gety{(b3)}{\G}
\draw[->] (n7-bot) |- ($(spx |- n7-bot)+(0,-3)$) -- (\SX,\G-5.6);
\pairstack{o3}{(\SX+5.8,\G-12.6)}{4}
\node[capt, text width=140pt] (c3) at (\SX+29,\G-5.5) {Programs (x86 asm) + annotations\\[-1pt]{\footnotesize\itshape One pair per branch mode $\times$ interleave variant}};
\path let \p1=(o3-ann-s), \p2=(c3.south) in coordinate (jn) at (\x1,{min(\y1-5,\y2-2.5)});

% ================= Stage 3 =================
\gety{(jn)}{\T}\pgfmathsetmacro\T{\T-7}
% gem5 container: logo + "processor model"; harness ring around the gray die
\draw[rounded corners=4pt, line width=0.6pt] (\IL,\T-15) rectangle (232,\T-85) coordinate (gemSE);
\node[font=\footnotesize\itshape, anchor=north east] (pm) at (228,\T-18.5) {processor model};
\node[anchor=east] at ($(pm.west)+(-1.5,0.2)$) {\includegraphics[height=6.5pt]{gem5logo.pdf}};
\draw[rounded corners=3pt] (68,\T-34) rectangle (226,\T-79);
\node[draw=sssrcink, rounded corners=2pt, fill=sssrcbg, accent=sssrcink, minimum width=146pt, minimum height=17pt, anchor=north west] (die) at (74,\T-39) {Defense implementation under test};
\foreach \x in {118,150,182} \draw[line width=0.5pt] (\x,\T-34) -- (\x,\T-39);
\foreach \y in {-43,-47.5,-52} \draw[line width=0.5pt] (220,\T+\y) -- (226,\T+\y);
\node[sm, anchor=north west, align=left] (hl) at (74,\T-60) {+\,Non-invasive harness: \textcolor{ssxmitink}{pipeline-stage probes};\\speculation and timing directives};
% inputs from the program + annotation pairs
\draw[->] (o3-asm-s) -- (o3-asm-s |- die.north);
\draw (o3-ann-s) -- (jn);
\draw[->] (jn) -- (jn |- 0,\T-34);
\node[dot] at (jn) {};
\node[q, anchor=north west, align=left] at ($(jn |- 0,\T-16.5)+(2,0)$) {speculation\\[-1.5pt]directives};
\draw[->] (jn) -- (236,0 |- jn) -- (236,\T-101);
\node[q, anchor=east] at (234.5,\T-94.5) {labels};
% pipeline info artifact -> inspector -> bug
\cardstack{tr}{(170,\T-103)}{1}{\gTrace}
\draw[->] (170,\T-79) -- (tr-n);
\node[sm, anchor=east, align=right] at ($(tr-w)+(-2.5,0)$) {Pipeline info\\artifact};
\node[box, anchor=west, minimum height=15pt, align=left, inner xsep=3pt] (I) at (194,\T-108.5) {\hspace*{8pt}Inspector\hspace*{3pt}};
\magn{($(I.west)+(6,0.8)$)}{2.5}
\draw[->] (tr-e) -- (tr-e -| I.west);
\draw[->, line width=0.7pt, ssxmitink] (I.south -| 214,0) -- (214,\T-122.5);
\buggl{(214,\T-131)}{1.6}{ssxmitink}
\node[font=\footnotesize\bfseries, text=ssxmitink, anchor=west] at (222.5,\T-131) {Bug?};
% sources: defense implementation -> die ; transmission mechanisms -> harness ("selects hooks") + inspector
\sot{gb}{(die.west -| lxc)}{\textless/\textgreater}
\node[sotl, text width=50.6pt] (gbl) at ($(gb-s)+(1.3,-1.5)$) {Defense implementation (gem5)};
\draw[->, sssrcink, line width=0.6pt] (gb-e) -- (die.west);
\draw[dashed, ->, sssrcink] ($(dsl.south)+(0,-3)$) -- ($(gb-n)+(0,3)$) coordinate[pos=0.5] (gm);
\node[q, rotate=90, anchor=south] at ($(gm)+(-2,0)$) {\textbf{SimSpect goal:} validate security guarantee on implementation};
\sot{tm}{(\LX,\T-104)}{}
\node[sotl] (tml) at ($(tm-s)+(0,-1.5)$) {Defense transmission mechanisms};

\draw[ssxmitink, line width=0.6pt] (tm-e) -- (61,0 |- tm-e) |- (200,\T-129);
\draw[->, ssxmitink, line width=0.6pt] (200,\T-129) -- (200,0 |- I.south);
\draw[->, ssxmitink, line width=0.6pt] (96,\T-129) -- (96,\T-79) node[q, text=ssxmitink, pos=0.2, anchor=west, xshift=1.5pt] {selects hooks};
\node[dot, ssxmitink] at (96,\T-129) {};
\draw[blob] (\PL,\T) rectangle (\PR,\T-139);
\node[btitle] (t3) at (\SX+20,\T) {\textbf{Stage 3}\enspace Execution};
\node[lang, anchor=west, fill=white, inner xsep=2pt, inner ysep=1pt] at ($(t3.east)+(0,0)$) {(gem5)};
\begin{pgfonlayer}{bg}\draw[sssrcbg, line width=6pt, line cap=round] (\LX,-53) -- (tm-n);\end{pgfonlayer}
\end{tikzpicture}
\end{document}
